\paragrafos{Εύρος, κορυφές και απόμακρες τιμές}{Γ.2.2, «Εύρος, κορυφές, πολλαπλές κορυφές, απόμακρες τιμές / Μεταβλητότητα δεδομένων»}

\begin{anagke}
Ο Νίκος και η Ρήνη παίζουν στο διάλειμμα ένα παιχνίδι με βολές στο καλάθι. Κάνουν πέντε γύρους των $10$ βολών ο καθένας. Ο Νίκος έβαλε $1$, $9$, $6$, $10$ και $4$ καλάθια. Η Ρήνη έβαλε $5$, $6$, $6$, $6$ και $7$.

«Ίδιος μέσος όρος!», φωνάζει ο Νίκος, «και οι δύο έχουμε $6$». Η Ρήνη κουνά το κεφάλι: «Εγώ ξέρω σχεδόν πάντα τι θα κάνω. Εσύ μια παίρνεις δέκα και μια ένα. Δεν παίζουμε το ίδιο». Η κυρία Αλεξάνδρα ρωτά: «Πώς θα έδειχνες με έναν αριθμό αυτό που βλέπει η Ρήνη;».

Η μέση τιμή δεν φτάνει για να περιγράψουμε δεδομένα. Χρειαζόμαστε και ένα μέτρο για το πόσο «απλώνονται» οι τιμές.

Σε αυτή την παράγραφο μαθαίνουμε το εύρος, την κορυφή (και τις πολλαπλές κορυφές), τις απόμακρες τιμές και τι σημαίνει μεταβλητότητα των δεδομένων.
\end{anagke}

\begin{lexi}{εύρος, κορυφή, απόμακρη τιμή, μεταβλητότητα}{ονόματα και σημασίες}
Το \lexiword{εύρος} είναι η έκταση (το «πλάτος») που καταλαμβάνουν τα δεδομένα. Η \lexiword{κορυφή} είναι η κορυφή ενός λόφου: η τιμή που «ξεπροβάλλει» στο διάγραμμα επειδή εμφανίζεται πιο συχνά (και λέγεται και \emph{επικρατούσα τιμή}). Η \lexiword{απόμακρη} τιμή βρίσκεται μακριά από τις άλλες. Η \lexiword{μεταβλητότητα} είναι το πόσο μεταβάλλονται, δηλαδή πόσο διαφέρουν μεταξύ τους, οι τιμές.
\end{lexi}

\begin{metafrasi}
Η μεγαλύτερη τιμή μείον τη μικρότερη. & $R=\text{μεγαλύτερη}-\text{μικρότερη}$\\
Η τιμή που εμφανίζεται τις περισσότερες φορές. & κορυφή (επικρατούσα τιμή)\\
Δύο ή περισσότερες τιμές με την ίδια, μέγιστη, συχνότητα. & πολλαπλές κορυφές\\
Μια τιμή που ξεχωρίζει πολύ από τις υπόλοιπες. & απόμακρη τιμή\\
Οι τιμές είναι πολύ διαφορετικές μεταξύ τους. & μεγάλη μεταβλητότητα\\
\end{metafrasi}

\tmima{Η θεωρία}

\begin{kanonas}{εύρος}
$$\text{εύρος }R=(\text{μεγαλύτερη τιμή})-(\text{μικρότερη τιμή})$$
Σε \textbf{ομαδοποιημένα} δεδομένα το εύρος είναι η διαφορά του μικρότερου άκρου της πρώτης κλάσης από το μεγαλύτερο άκρο της τελευταίας κλάσης.
\end{kanonas}

Στο παιχνίδι του Νίκου και της Ρήνης, τα δεδομένα στην αριθμογραμμή δείχνουν καθαρά τη διαφορά.

\begin{center}
\begin{tikzpicture}[x=0.8cm]
\draw[-{Stealth},thick] (-0.5,0)--(11,0);
\foreach \x in {0,...,10}{\draw (\x,0.1)--(\x,-0.1) node[below,font=\scriptsize]{$\x$};}
\foreach \x in {1,4,6,9,10}{\fill[gwarmD] (\x,0.5) circle (3.5pt);}
\foreach \x in {5,7}{\fill[gblue] (\x,1.2) circle (3.5pt);}
\fill[gblue] (6,1.2) circle (3.5pt);\fill[gblue] (6,1.55) circle (3.5pt);\fill[gblue] (6,1.9) circle (3.5pt);
\node[font=\small,gwarmD,anchor=west] at (11.2,0.5) {Νίκος: $R=9$};
\node[font=\small,gblue,anchor=west] at (11.2,1.4) {Ρήνη: $R=2$};
\end{tikzpicture}
\end{center}

\begin{kanonas}{κορυφή}
\textbf{Κορυφή} (ή \textbf{επικρατούσα τιμή}) είναι η τιμή που εμφανίζεται τις περισσότερες φορές, δηλαδή έχει τη μεγαλύτερη συχνότητα. Αν δύο ή περισσότερες τιμές έχουν την ίδια μεγαλύτερη συχνότητα, έχουμε \textbf{πολλαπλές κορυφές}.
\end{kanonas}

Η κορυφή έχει νόημα και σε δεδομένα που δεν είναι αριθμοί. Αν ρωτήσουμε ποιο είναι το αγαπημένο χρώμα, η κορυφή μπορεί να είναι «μπλε», ενώ μέση τιμή χρωμάτων δεν υπάρχει. Σε ιστόγραμμα κορυφή είναι η κλάση με το ψηλότερο ορθογώνιο. Το ιστόγραμμα εδώ έχει δύο κορυφές:

\begin{center}
\binhist{2,7,3,6,1}{0,10,20,30,40,50}{8}{διάρκεια διαδρομής (λεπτά)}
\end{center}

\begin{orismos}{απόμακρη τιμή}
\textbf{Απόμακρη τιμή} λέγεται το δεδομένο που διαφέρει σημαντικά από τα υπόλοιπα δεδομένα.
\end{orismos}

Πού οφείλονται οι απόμακρες τιμές; Συχνά σε \emph{λάθος}: κάποιος πληκτρολόγησε $85$ αντί για $8{,}5$ ή το όργανο μέτρησης δεν δούλευε σωστά. Άλλες φορές είναι \emph{φυσική ποικιλία}: ανάμεσα σε παιδιά σε μια παιδική χαρά, ο παππούς είναι «απόμακρος» στην ηλικία, χωρίς να υπάρχει λάθος. Δεν υπάρχει αριθμητικός τύπος που να μας λέει πότε μια τιμή είναι απόμακρη: κοιτάμε τα δεδομένα και ρωτάμε αν μια τιμή «ξεκόβει» από τις άλλες.

Η απόμακρη τιμή μεγαλώνει πολύ το εύρος και τραβά τη μέση τιμή, ενώ η διάμεσος και η κορυφή αλλάζουν λίγο ή καθόλου.

\begin{idiotita}{μεταβλητότητα}
Όσο μεγαλύτερο είναι το εύρος, τόσο περισσότερο απλώνονται τα δεδομένα και τόσο μεγαλύτερη είναι η \textbf{μεταβλητότητά} τους. Για να συγκρίνουμε δύο σύνολα δεδομένων κοιτάμε μαζί τη μέση τιμή (πού βρίσκονται) και το εύρος (πόσο απλώνονται). Δύο σύνολα μπορεί να έχουν την ίδια μέση τιμή και πολύ διαφορετικό εύρος.
\end{idiotita}

\begin{prosoxi}
\begin{itemize}[leftmargin=1.2em,itemsep=0.3ex,topsep=2pt]
\item Η κορυφή είναι η \emph{τιμή}, όχι η συχνότητά της. Στα δεδομένα $2,5,5,5,5,7$ η κορυφή είναι το $5$ και όχι το $4$ (το $4$ είναι πόσες φορές εμφανίζεται).
\item Το εύρος είναι αριθμός που δεν είναι ποτέ αρνητικός. Είναι πάντα μεγαλύτερο μείον μικρότερο, και ας είναι οι τιμές αρνητικές: για $-3$ και $4$ βγαίνει $4-(-3)=7$.
\item Δεν πετάμε μια απόμακρη τιμή χωρίς λόγο. Πρώτα ρωτάμε από πού προέρχεται, αν είναι λάθος ή αληθινή.
\item Το εύρος κοιτά μόνο τα δύο άκρα. Μια απόμακρη τιμή αρκεί για να το μεγαλώσει πολύ.
\end{itemize}
\end{prosoxi}

\tmima{Λυμένα παραδείγματα}

\paradeigma{Το εύρος}
\tixero Η Πολυξένη κατέγραψε στη Θεσσαλονίκη τις θερμοκρασίες επτά ημερών ($^\circ$C): $-3$, $0$, $2$, $-1$, $4$, $1$, $2$.\par
\tithelo Το εύρος.\par
\possindeo Βρίσκω τη μεγαλύτερη και τη μικρότερη τιμή.\par
\apantisi Μεγαλύτερη $4$, μικρότερη $-3$. Εύρος $R=4-(-3)=7$ $^\circ$C.\par
\elegxos Από το $-3$ ως το $4$ στην αριθμογραμμή υπάρχουν $7$ μονάδες.

\paradeigma{Η κορυφή και οι πολλαπλές κορυφές}
\tixero Τα νούμερα παπουτσιών $12$ παιδιών: $36$, $37$, $37$, $38$, $36$, $39$, $37$, $38$, $36$, $40$, $38$, $37$.\par
\tithelo Την κορυφή.\par
\possindeo Μετρώ πόσες φορές εμφανίζεται κάθε τιμή.\par
\apantisi Το $36$ εμφανίζεται $3$ φορές, το $37$ εμφανίζεται $4$ φορές, το $38$ εμφανίζεται $3$ φορές, τα $39$ και $40$ από μία. Κορυφή το $37$.\par
\elegxos Οι συχνότητες αθροίζουν $3+4+3+1+1=12$, όσα και τα παιδιά. Στα δεδομένα $5,6,6,7,8,8,9$ τα $6$ και $8$ εμφανίζονται από δύο φορές και τα άλλα από μία: υπάρχουν δύο κορυφές, το $6$ και το $8$.

\paradeigma{Απόμακρη τιμή}
\tixero Σε ένα οικογενειακό τραπέζι οι ηλικίες των παιδιών είναι $9, 10, 11, 11, 12, 12, 12$ και ο παππούς Γιάννης είναι $83$.\par
\tithelo Να δούμε πώς επηρεάζει ο παππούς τη μέση τιμή, τη διάμεσο και το εύρος.\par
\possindeo Υπολογίζω πρώτα χωρίς τον παππού και μετά με τον παππού.\par
\apantisi Χωρίς τον παππού: άθροισμα $77$, μέση τιμή $77:7=11$, διάμεσος $11$, εύρος $12-9=3$. Με τον παππού: άθροισμα $160$, μέση τιμή $160:8=20$, διάμεσος $\dfrac{11+12}{2}=11{,}5$, εύρος $83-9=74$.\par
\elegxos Η μέση τιμή σχεδόν διπλασιάστηκε και το εύρος έγινε $74$, ενώ η διάμεσος αλλάζει ελάχιστα. Η μέση ηλικία $20$ δεν περιγράφει κανέναν στο τραπέζι. Η $83$ είναι απόμακρη, αλλά όχι λάθος: είναι ο παππούς.

\paradeigma{Ίδιος μέσος όρος, διαφορετική μεταβλητότητα}
\tixero Βολές του Νίκου: $1,9,6,10,4$. Βολές της Ρήνης: $5,6,6,6,7$.\par
\tithelo Να συγκρίνουμε τους δύο παίκτες.\par
\possindeo Υπολογίζω μέση τιμή και εύρος για τον καθένα.\par
\apantisi Νίκος: μέση τιμή $\frac{30}{5}=6$, εύρος $10-1=9$. Ρήνη: μέση τιμή $\frac{30}{5}=6$, εύρος $7-5=2$. Και οι δύο έχουν $6$ κατά μέσο όρο, αλλά η Ρήνη είναι πολύ πιο σταθερή (μικρό εύρος), ενώ ο Νίκος είναι απρόβλεπτος (μεγάλο εύρος).\par
\elegxos Το σύνολο της Ρήνης είναι «μαζεμένο» γύρω από το $6$ και του Νίκου απλωμένο από το $1$ ως το $10$, όπως φαίνεται στην αριθμογραμμή.

\paradeigma{Εύρος σε ομαδοποιημένα δεδομένα}
\tixero Τα ύψη $20$ παιδιών (cm): $140$--$145$: $2$ παιδιά, $145$--$150$: $5$, $150$--$155$: $8$, $155$--$160$: $4$, $160$--$165$: $1$.\par
\tithelo Το εύρος και την κλάση με τη μεγαλύτερη συχνότητα.\par
\possindeo Εύρος: μεγαλύτερο άκρο της τελευταίας κλάσης μείον μικρότερο άκρο της πρώτης.\par
\apantisi Εύρος $165-140=25$ cm. Το ψηλότερο ορθογώνιο είναι της κλάσης $150$--$155$ ($8$ παιδιά).\par
\elegxos Το εύρος είναι προσεγγιστικό, γιατί δεν ξέρουμε το πραγματικό ύψος του ψηλότερου ή του πιο κοντού παιδιού.

\tmima{Ασκήσεις}

\zoni{Το κατάλαβα}{εύρος, κορυφή, απόμακρη τιμή}

\ex Βρες το εύρος: α) $3,8,5,12,7$. β) $-4,0,6,-2$. γ) $15,15,15$. δ) $0{,}5$, $2{,}5$, $1{,}75$.

\ex Βρες την κορυφή ή τις κορυφές: α) $2,3,3,5,3,7$. β) $1,4,4,6,6,9$. γ) κόκκινο, μπλε, μπλε, πράσινο, κόκκινο, μπλε.

\ex Σωστό ή λάθος; α) Το εύρος μπορεί να είναι αρνητικός αριθμός. β) Αν το εύρος είναι $0$, όλες οι τιμές είναι ίσες. γ) Η κορυφή είναι πάντα μία. δ) Η κορυφή είναι η τιμή που εμφανίζεται τις περισσότερες φορές. ε) Μια απόμακρη τιμή μεγαλώνει το εύρος.

\ex Βρες, αν υπάρχει, την απόμακρη τιμή: α) $4,5,6,5,4,48$. β) $120,118,125,3,122$. γ) $10,12,11,13,12$.

\ex Το ιστόγραμμα δείχνει πόσοι μαθητές της Α3 μένουν σε κάθε απόσταση από το σχολείο (σε λεπτά με τα πόδια).
\begin{center}
\binhist{2,7,3,6,1}{0,10,20,30,40,50}{8}{χρόνος (λεπτά)}
\end{center}
α) Πόσοι μαθητές ρωτήθηκαν; β) Ποιο είναι το εύρος; γ) Έχει το ιστόγραμμα περισσότερες από μία κορυφές; δ) Πόσοι μαθητές χρειάζονται λιγότερο από $20$ λεπτά;

\ex Οι θερμοκρασίες μιας εβδομάδας στο χωριό ($^\circ$C), από Δευτέρα ως Κυριακή: $-2, 3, 0, 5, 2, 8, -2$. Βρες τη μέση τιμή, το εύρος και την κορυφή.

\ex Δύο σύνολα δεδομένων: Α: $5,6,7,8,9$ και Β: $1,4,7,10,13$. Βρες τη μέση τιμή και το εύρος του καθενός. Ποιο έχει μεγαλύτερη μεταβλητότητα;

\zoni{Το χρησιμοποιώ}{διαγωνισμοί, πωλήσεις, διαφημίσεις}

\ex \textbf{Οι δρομείς.} Σε πέντε αγώνες των $100$ m η Αγγελική έκανε $16, 17, 17, 17, 18$ δευτερόλεπτα και η Πολυξένη $14, 15, 17, 19, 20$. Βρες τη μέση τιμή και το εύρος της καθεμιάς. Ποια είναι πιο σταθερή; Ποια έκανε τον καλύτερο χρόνο;

\ex \textbf{Τα παγωτά.} Σε ένα περίπτερο στην παραλία πουλήθηκαν επτά μέρες $40, 36, 44, 38, 42, 40$ και $5$ παγωτά. Την τελευταία μέρα έβρεχε. α) Βρες μέση τιμή, διάμεσο και εύρος και με τις επτά μέρες και χωρίς τη βροχερή. β) Ποιο μέτρο άλλαξε λιγότερο;

\ex \textbf{Τα ύψη.} Στον πίνακα φαίνονται τα ύψη $20$ παιδιών:\par
\begin{center}
\renewcommand{\arraystretch}{1.25}
\begin{tabular}{lccccc}
\toprule
Ύψος (cm) & $140$--$145$ & $145$--$150$ & $150$--$155$ & $155$--$160$ & $160$--$165$\\
Παιδιά & $2$ & $5$ & $8$ & $4$ & $1$\\
\bottomrule
\end{tabular}
\end{center}
Βρες το εύρος και τη μέση τιμή του ύψους (με τα κέντρα των κλάσεων).

\ex \textbf{Οι βαθμοί.} Ο πίνακας δείχνει τους βαθμούς $22$ μαθητών σε ένα τεστ (κλίμακα $0$--$20$):\par
\begin{center}
\renewcommand{\arraystretch}{1.25}
\begin{tabular}{lcccccc}
\toprule
Βαθμός & $12$ & $13$ & $14$ & $15$ & $16$ & $17$\\
Μαθητές & $2$ & $3$ & $6$ & $6$ & $3$ & $2$\\
\bottomrule
\end{tabular}
\end{center}
Βρες την κορυφή (ή τις κορυφές), το εύρος και τη μέση τιμή.

\ex \textbf{Το ζάρι.} Ο Ηλίας έριξε ένα ζάρι $12$ φορές και έγραψε: $3,5,2,6,5,1,5,4,6,5,2,3$. Φτιάξε πίνακα συχνοτήτων και βρες την κορυφή, το εύρος και τη μέση τιμή.

\zoni{Το σκέφτομαι}{τι αλλάζει και τι μένει}

\ex Βρες πέντε αριθμούς με εύρος $10$, διάμεσο $7$ και κορυφή $7$.

\ex Τα δεδομένα $2, 5, 8$ έχουν μέση τιμή $5$ και εύρος $6$. α) Πρόσθεσε $3$ σε κάθε τιμή. Τι γίνεται η μέση τιμή και τι το εύρος; β) Πολλαπλασίασε κάθε τιμή με $2$. Τι γίνεται τώρα; Τι παρατηρείς;

\ex Τα δεδομένα $3,4,5,5,6,7$ έχουν μέση τιμή $5$ και διάμεσο $5$. Πρόσθεσε το $40$. Βρες τη νέα μέση τιμή και τη νέα διάμεσο. Τι κάνει η απόμακρη τιμή σε καθένα από τα δύο;

\ex Στο σχολικό πανηγύρι ρωτήσαμε τις ηλικίες όσων ήρθαν.
\begin{center}
\binhist{2,7,1,3,6}{0,10,20,30,40,50}{8}{ηλικία (έτη)}
\end{center}
α) Πόσες κορυφές έχει το ιστόγραμμα; β) Πώς το εξηγείς; γ) Βρες τη μέση τιμή (με τα κέντρα των κλάσεων). Περιγράφει κάποιον από τους παρευρισκόμενους;

\zoni{Για το σπίτι}{μέτρησε και σύγκρινε}

\ex Μέτρησε επί πέντε μέρες πόσα λεπτά κάνεις να φτάσεις στο σχολείο. Βρες τη μέση τιμή, τη διάμεσο και το εύρος.

\ex Ρώτησε $12$ ανθρώπους ποιο είναι το αγαπημένο τους χρώμα. Βρες την κορυφή. Γιατί δεν έχει νόημα να βρεις μέση τιμή;
